% © 2026 Ing. David Jaroš — CC BY-NC-ND 4.0
%
% This work is licensed under a Creative Commons Attribution-NonCommercial-NoDerivatives
% 4.0 International License (CC BY-NC-ND 4.0).
%
% License History: Earlier drafts (up to v0.3) were released under CC BY 4.0.
% From v0.4 onward, all material is released under CC BY-NC-ND 4.0 to protect
% the integrity of the theoretical work during ongoing academic development.
%
% See LICENSE.md for full license text.
%
% RETRACTED 2026-07-17: this file conflated an FPE for P with a diffusion
% equation for Theta. It is retained as a transparent historical audit record.

% UBT-AI-PROVENANCE-BEGIN
% schema: ubt-ai-provenance/v1
% tier: B_machine_verified
% ai_assistance: disclosed
% human_review: machine-verification
% editorial_responsibility: Ing. David Jaroš
% policy: ../../AI_PROVENANCE.md
% notice: Machine-verified against named sources or verifiers; individual attestation is not claimed.
% UBT-AI-PROVENANCE-END

\ifdefined\INCLUDEMODE
\else
  \documentclass[12pt,a4paper]{article}
  \usepackage[a4paper, margin=2.5cm]{geometry}
  \usepackage{amsmath, amssymb, amsthm}
  \usepackage{mathtools}
  \usepackage{hyperref}
  \usepackage{xcolor}
  \usepackage{mdframed}
  \usepackage{booktabs}

  \newtheorem{theorem}{Theorem}[section]
  \newtheorem{lemma}[theorem]{Lemma}
  \newtheorem{proposition}[theorem]{Proposition}
  \newtheorem{corollary}[theorem]{Corollary}
  \theoremstyle{definition}
  \newtheorem{definition}[theorem]{Definition}
  \theoremstyle{remark}
  \newtheorem{remark}[theorem]{Remark}

  \newmdenv[backgroundcolor=yellow!20, linecolor=orange, linewidth=1.5pt]{gapbox}
  \newmdenv[backgroundcolor=green!15, linecolor=green!60!black, linewidth=1.5pt]{resultbox}
  \newmdenv[backgroundcolor=blue!8, linecolor=blue!50, linewidth=1.5pt]{insightbox}

  \title{\textbf{Step 7: Retracted Born-Rule Claim}\\
         \large Historical audit of norm decay under amplitude diffusion}
  \author{David Jaro\v{s}}
  \date{2026-03-06}

  \begin{document}
  \maketitle

  \begin{abstract}
  This historical file previously claimed that
  $P=\mathrm{Sc}(\Theta^\dagger\Theta)$ is automatically conserved when
  $\Theta$ satisfies a diffusion/Fokker--Planck-type equation.  That claim is
  retracted.  The calculation retained below instead shows a negative norm
  derivative for a genuine diffusion term.  The Born rule and unitary physical
  evolution therefore remain open in \texttt{CLAIMS\_MATRIX.md}.
  \end{abstract}
\fi

\begin{center}
\fbox{\begin{minipage}{0.92\textwidth}
\textbf{RETRACTED (2026-07-17).}
The previous argument conflated the Fokker--Planck equation for a probability
density $P$ with a diffusion equation for the amplitude $\Theta$.  For example,
$\Theta=e^{ikx-Dk^2T}$ solves $\partial_T\Theta=D\partial_x^2\Theta$, while
$\int|\Theta|^2dx$ decays as $e^{-2Dk^2T}$.  This file is retained to preserve
the correct norm-decay computation and the audit trail.  It is superseded by
the OPEN Born-rule/unitarity status in \texttt{CLAIMS\_MATRIX.md}.
\end{minipage}}
\end{center}

%──────────────────────────────────────────────────────────────────────────────
\section{Setup: Norm in $\mathbb{C}\otimes\mathbb{H}$}
\label{sec:setup}
%──────────────────────────────────────────────────────────────────────────────

For $\Theta \in \mathcal{B} = \mathbb{C}\otimes\mathbb{H}$, the natural
norm is:
\begin{equation}
  |\Theta|^2 \;:=\; \mathrm{Sc}[\Theta^\dagger\Theta]
             \;=\; \frac{1}{2}\mathrm{Tr}[\Theta^\dagger\Theta]
  \label{eq:norm}
\end{equation}
where $\Theta^\dagger$ is the biquaternionic conjugate and $\mathrm{Sc}$
denotes the scalar part.  This is:
\begin{enumerate}
  \item \emph{Real}: $\mathrm{Sc}[\Theta^\dagger\Theta] \in \mathbb{R}_{\geq 0}$.
  \item \emph{Positive definite}: $|\Theta|^2 = 0 \Leftrightarrow \Theta = 0$.
  \item \emph{Gauge invariant}: $|g\Theta|^2 = |\Theta|^2$ for $g \in SU(2)$ or $U(1)$.
\end{enumerate}
The probability density candidate is $P(x,\psi) = |\Theta(x,\psi)|^2$.

%──────────────────────────────────────────────────────────────────────────────
\section{Retracted conservation claim and audit calculation}
\label{sec:born}
%──────────────────────────────────────────────────────────────────────────────

The former claim was that, if $\Theta(T,Q)$ obeys
\begin{equation}
  \partial_T\Theta=-\nabla_Q\!\cdot(A\Theta)+D\nabla_Q^2\Theta,
  \label{eq:FPE-amplitude}
\end{equation}
then $P=\mathrm{Sc}(\Theta^\dagger\Theta)$ obeys the corresponding
Fokker--Planck continuity equation.  The following calculation shows why that
implication fails.

Using
\begin{align}
  \partial_T\,\mathrm{Sc}(\Theta^\dagger\Theta)
  &=\mathrm{Sc}\!\left[(\partial_T\Theta^\dagger)\Theta
       +\Theta^\dagger(\partial_T\Theta)\right],
\end{align}
the diffusion term satisfies
\begin{align}
  \mathrm{Sc}\!\left[(\nabla_Q^2\Theta)^\dagger\Theta
       +\Theta^\dagger\nabla_Q^2\Theta\right]
  &=\nabla_Q^2\,\mathrm{Sc}(\Theta^\dagger\Theta)
    -2\,\mathrm{Sc}\!\left[(\nabla_Q\Theta)^\dagger\nabla_Q\Theta\right].
\end{align}
After integration and vanishing boundary flux, the second term remains:
\begin{equation}
  \partial_T\int P\,dQ
  =-2D\int\mathrm{Sc}\!\left[(\nabla_Q\Theta)^\dagger
      \nabla_Q\Theta\right]dQ\leq0,
  \label{eq:norm-decay}
\end{equation}
up to any separately analysed drift contribution.  Thus amplitude diffusion is
contractive in norm, not unitary in general.

\paragraph{Explicit counterexample.}
On a periodic interval of length $L$, let
\begin{equation}
  \Theta(x,T)=e^{ikx-Dk^2T}.
\end{equation}
It solves $\partial_T\Theta=D\partial_x^2\Theta$, but
\begin{equation}
  \int_0^L|\Theta(x,T)|^2dx=L e^{-2Dk^2T},
\end{equation}
which is nonconstant for $k\neq0$ and $D>0$.

A standard Fokker--Planck equation written directly for a density $P$,
$\partial_TP=-\nabla\cdot(AP)+D\nabla^2P$, conserves $\int P$ under suitable
boundary conditions.  That fact does not imply that
$P=\mathrm{Sc}(\Theta^\dagger\Theta)$ obeys the same equation when $\Theta$
obeys \eqref{eq:FPE-amplitude}.

%──────────────────────────────────────────────────────────────────────────────
\section{Current status}
\label{sec:status}
%──────────────────────────────────────────────────────────────────────────────

\begin{gapbox}
\textbf{Gap S2 / Born rule / unitarity: OPEN.}

The positive scalar $\mathrm{Sc}(\Theta^\dagger\Theta)$ remains a natural
candidate density, but its Born interpretation and conservation require a
separate derivation from the finalized canonical UBT dynamics.  In particular,
one must identify a norm-preserving physical evolution, or demonstrate how a
reversible higher-layer dynamics yields an effective probabilistic equation
without losing the required inner-product structure.

The previous ``Gap S2: CLOSED'' label and all claims of automatic FPE norm
conservation are withdrawn.
\end{gapbox}

\begin{center}
\small
\begin{tabular}{@{}p{0.34\textwidth}p{0.19\textwidth}p{0.37\textwidth}@{}}
  \toprule
  Result & Current status & Note \\
  \midrule
  Positivity of $\mathrm{Sc}(\Theta^\dagger\Theta)$ & Candidate property & depends on chosen involution/representation \\
  Amplitude-diffusion norm conservation & Retracted & contradicted by \eqref{eq:norm-decay} \\
  Born rule from canonical UBT & Open gap & see \texttt{CLAIMS\_MATRIX.md} \\
  Unitary physical evolution & Open gap & must be derived \\
  \bottomrule
\end{tabular}
\end{center}

\ifdefined\INCLUDEMODE
\else
  \end{document}
\fi
